\renewcommand{\Anno}{2002}

\renewcommand{\NumeroProblema}{4}

\renewcommand{\Testo}{Un utente di internet, Paolo, vuole scaricare un
documento dalla rete.  Per connettersi ha due possibilit\`{a}.
\begin{itemize}
	\item Un collegamento mediante una linea telefonica $A$ ad un
	server $S_{A}$: in questo caso la linea telefonica $A$ cadr\`{a}
	dopo un numero intero $n_{A}$ di minuti, $1\leq n_{A}\leq 25$,
	tutti equiprobabili, mentre il server $S_{A}$ si pu\`{o} spegnere
	dopo un intervallo di tempo di 1, 2, \ldots, 20 minuti, ognuno con
	uguale probabilit\`{a}, indipendentemente dalla durata della
	telefonata.
	
	\item Un collegamento mediante una linea telefonica $B$ ad un
	server $S_{B}$: in questo caso la linea telefonica $B$ cadr\`{a}
	dopo un numero intero $n_{B}$ di minuti, $1\leq n_{B}\leq 100$,
	tutti equiprobabili, mentre il server $S_{B}$ si pu\`{o} spegnere
	dopo un intervallo di tempo di 1, 2, \ldots, 15 minuti, ognuno con
	uguale probabilit\`{a}, indipendentemente dalla durata della
	telefonata.
\end{itemize}
	
Per ottenere l'intero documento \`{e} necessario rimanere collegati
attraverso una linea telefonica ad un server senza interruzioni per
almeno 6 minuti.
\begin{enumerate}
	\item A quale linea conviene collegarsi per ottenere il documento
	al primo tentativo?
	
	\item Paolo non \`{e} interessato ad ottenere il documento al
	primo tentativo ma ha deciso di ripetere pi\`{u} tentativi di
	collegamento con la stessa linea fino ad arrivare ad una
	probabilit\`{a} di successo del 95\%.
		
	Determinare il numero di tentativi necessari per ciascuna delle
	due linee telefoniche.
	
	\item Paolo deve scaricare un elevato numero di documenti, tutti
	della durata di 6 minuti.  Ne scarica uno al giorno con il
	seguente procedimento: se entro 3 tentativi non riesce ad ottenere
	il documento, allora se lo fa spedire per corriere al costo di 5
	euro.  Alla lunga, sapendo che la linea $A$ costa 25 centesimi di
	euro per ogni chiamata, mentre la linea $B$ costa 30 centesimi per
	ogni chiamata, determinare la spesa media per documento che Paolo
	dovr\`{a} sostenere usando sempre la linea $A$ o sempre la linea
	$B$.
\end{enumerate}}

\renewcommand{\Risposta}{\begin{enumerate}
		\item  Conviene la linea $B$.
	
		\item  Con la linea $A$ il numero di tentativi \`{e} 4, con la 
		linea $B$ \`{e} 3.
	
		\item Con la linea $A$ la spesa media a documento \`{e} 73
		centesimi, con la linea $B$ \`{e} $9407/135$ centesimi (circa
		69,68 centesimi).
	\end{enumerate}}

\renewcommand{\Aiutino}{\begin{enumerate}
	\renewcommand{\labelenumi}{(\alph{enumi})} 
	\item Qual \`{e} la probabilit\`{a} che la linea telefonica $A$
	``sopravviva'' almeno 6 minuti?  E che il server $S_{A}$
	``sopravviva'' almeno 6 minuti?  Idem per la linea $B$ ed il
	server $S_{B}$.

	\item Qual \`{e} la probabilit\`{a} di fallire una volta con la
	linea $A$?  E di fallire $n$ volte?  Idem per la linea $B$.  Non
	resta che trovare $n$ in modo che tali probabilit\`{a} siamo
	minori di $1/20$.

	\item  Per determinare la spedia media bisogna calcolare il costo di un 
	collegamento per la probabilit\`{a} di farcela in un tentativo, pi\`{u} 
	il costo di 2 collegamenti per la probabilit\`{a} di farcela in 2 
	tentativi, pi\`{u} \ldots
\end{enumerate}}

\renewcommand{\SoluzioneUfficiale}{\TitoloSoluzione \DomandaA Per 
ottenere il documento \`{e} necessario che sia la linea telefonica 
sia il server funzionino senza interruzioni per almeno 6 minuti.

Per quanto riguarda la linea telefonica $A$ abbiamo quindi $25-5=20$
casi favorevoli (ovvero i casi in cui la telefonata dura 6, 7, \ldots,
25 minuti) su 25 casi totali, mentre per il server $S_{A}$ abbiamo
$20-5=15$ casi favorevoli (ovvero i casi in cui il server si spegne
dopo 6, 7, \ldots, 20 minuti) sui 20 totali.  Avendo per ipotesi che i
2 eventi sono indipendenti, se indichiamo con $P_{A}$ la
probabilit\`{a} di successo con la linea $A$ si ha quindi
$$P_{A}=\frac{20}{25}\cdot\frac{15}{20}=\frac{3}{5}.$$

Analogamente con la linea telefonica $B$ abbiamo $100-5=95$
casi favorevoli sui 100 totali, mentre per il server $S_{B}$ abbiamo
$15-5=10$ casi favorevoli sui 15 totali.  Avendo per ipotesi che i 2
eventi sono indipendenti, se indichiamo con $P_{B}$ la probabilit\`{a}
di successo con la linea $B$ si ha quindi
$$P_{B}=\frac{95}{100}\cdot\frac{10}{15}=\frac{19}{30}.$$

La probabilit\`{a} di successo al primo tentativo \`{e} dunque 
maggiore con la linea $B$.

\DomandaB
Indichiamo con $I_{A}$ e $I_{B}$ le probabilit\`{a} di insuccesso con 
la linea $A$ e la linea $B$, rispettivamente, date da
$$I_{A}=1-P_{A}=\frac{2}{5},
\hspace{4em}
I_{B}=1-P_{B}=\frac{11}{30}.$$

Le probabilit\`{a} di insuccesso in $n$ tentativi con le due linee
saranno quindi, rispettivamente, $(I_{A})^{n}$ e $(I_{B})^{n}$. Ora la 
domanda del testo \`{e} equivalente a determinare, per ciascuno dei 
due casi, il pi\`{u} piccolo $n$ per cui tali espressioni sono minori 
del 5\%, cio\`{e} di $1/20$.

Nel caso della linea $A$ il pi\`{u} piccolo $n$ \`{e} 4, in quanto
$$\left(\frac{2}{5}\right)^{3}=\frac{8}{125}>\frac{1}{20},
\hspace{4em}
\left(\frac{2}{5}\right)^{4}=\frac{16}{625}<\frac{1}{20}.$$

Per la linea $B$ il pi\`{u} piccolo $n$ \`{e} 3, in quanto
$$\left(\frac{11}{30}\right)^{2}>\left(\frac{1}{3}\right)^{2}>\frac{1}{20},
\hspace{4em}
\left(\frac{11}{30}\right)^{3}=\frac{1331}{27\cdot 50\cdot 20}=
\frac{1331}{1350}\cdot\frac{1}{20}<\frac{1}{20}.$$

\DomandaC
Con la linea $A$ abbiamo probabilit\`{a} $P_{A}$ di farcela in un 
tentativo (con una spesa di 25 centesimi), probabilit\`{a} 
$I_{A}\cdot P_{A}$ di farcela al secondo tentativo (con una spesa di 
50 centesimi), probabilit\`{a} $(I_{A})^{2}\cdot P_{A}$ di farcela al 
terzo (con una spesa di 75 centesimi), probabilit\`{a} $(I_{A})^{3}$ 
di non farcela in 3 tentativi e dover ricorrere al corriere (con una 
spesa di 75 centesimi per il telefono e 500 centesimi per il 
corriere). Il costo medio sar\`{a} pertanto
\begin{eqnarray*}
	C_{A} & = & P_{A}\cdot 25+I_{A}\cdot P_{A}\cdot 50+I_{A}^{2}\cdot
	P_{A}\cdot 75+I_{A}^{3}\cdot 575 \\
	& = & \frac{3}{5}\cdot 25+\frac{2}{5}\cdot \frac{3}{5}\cdot
	50+\left(\frac{2}{5}\right)^{2}\cdot \frac{3}{5}\cdot
	75+\left(\frac{2}{5}\right)^{3}\cdot 575=73
\end{eqnarray*}
centesimi di euro.

Ragionando analogamente per la linea $B$, e tenendo conto che ora il 
costo di ogni telefonata \`{a} di 30 centesimi, avremo che il costo 
medio \`{e}
\begin{eqnarray*}
	C_{B} & = & P_{B}\cdot 30+I_{B}\cdot P_{B}\cdot 60+I_{B}^{2}\cdot
	P_{B}\cdot 90+I_{B}^{3}\cdot 590 \\
	& = & \frac{19}{30}\cdot 30+\frac{11}{30}\cdot \frac{19}{30}\cdot
	60+\left(\frac{11}{30}\right)^{2}\cdot \frac{19}{30}\cdot
	90+\left(\frac{11}{30}\right)^{3}\cdot 590=\frac{9407}{135},
\end{eqnarray*}
pari a circa $69.68$ centesimi di euro. }

\renewcommand{\SoluzioniAlternative}{
}

\renewcommand{\Osservazioni}{
}

\renewcommand{\Commenti}{
}